\chapter{Base teórica}
\label{cha:base-teorica}

\section{Introducción}
\label{sec:intro-base-teorica}

Antes de entrar en el diseño concreto del sistema, conviene fijar el terreno
teórico sobre el que se apoya: qué se entiende por arquitectura de
supervisión industrial, qué papel juega N3uron dentro de ella, cómo
funcionan los LLM y por qué el Model Context Protocol resulta la pieza que
permite conectarlos. El capítulo empieza revisando qué se ha hecho ya en la
integración de modelos de lenguaje con sistemas industriales, para situar
frente a ese panorama las decisiones tomadas en este TFG.

\section{Estado del arte}
\label{sec:estado-arte}

La integración de modelos de lenguaje de gran tamaño en el ámbito de la automatización
industrial es una línea de trabajo extremadamente novedosa, que lleva muy pocos años en
desarrollo, y esto se refleja en los resultados: la mayoría de las propuestas hechas
alrededor de este tema quedan en un simple prototipo, sin llegar a mayor desarrollo. Los
resultados de una revisión sistemática publicada en mayo de 2026, con un corpus de 88
publicaciones seleccionadas entre un bloque de más de 2300, confirman lo expuesto con
cifras concretas: el 75~\% de los sistemas descritos se sitúa en un nivel de madurez
tecnológica bajo (entre el prototipo y la validación temprana), y solo el 9,1~\%
evidencia haber llegado a introducirse en producción \cite{henkel2026}. En cuanto a los
usos para los que se han diseñado estos sistemas, los objetivos más habituales son la
asistencia al usuario, la supervisión de procesos y la optimización, sin que predomine
el control directo de la producción.

Sobre este panorama, existen diversos trabajos que abordan un problema común: un LLM
aislado es incapaz de conocer el estado del sistema, por lo que es necesario un
mecanismo que le dote de contexto sobre la instalación. Xia et al. proponen un modelo de
información basado en eventos, diseñado específicamente para alimentar al modelo con
datos del sistema de automatización en distintos niveles \cite{xia2024}. Vyas y
Mercangöz, por su parte, plantean una arquitectura de validación y reformulación apoyada
en un gemelo digital: antes de ejecutar cualquier acción, el agente la valida sobre la
simulación, y si se detecta un fallo, otro agente reformula la petición
\cite{vyas-mercangoz}.

Se aprecia, sin embargo, un rasgo común entre estas dos propuestas que resulta muy
relevante para este TFG: la conexión entre el LLM y el sistema industrial se resuelve
mediante un mecanismo construido a medida para cada proyecto. Ninguno de los dos emplea
un protocolo estandarizado para la comunicación entre modelo y sistema. Es justamente el
Model Context Protocol, presentado en el apartado~\ref{sec:model-context-protocol}, el
encargado de resolver este problema, ofreciendo un mecanismo genérico y reutilizable
para la conexión del LLM con cualquier fuente de datos o herramienta externa.

Además, el enfoque entre los trabajos mencionados y el presente TFG es ligeramente
distinto. Mientras que Xia et al. y Vyas y Mercangöz orientan sus sistemas al control del
proceso en tiempo real, este trabajo se centra en el análisis de datos históricos —en
este caso, de consumo eléctrico— con el fin de servir como herramienta de apoyo a la
toma de decisiones, sin suponer una intervención directa sobre el proceso. Esta
diferencia de enfoque se traduce en un agente con un perfil de riesgo mucho más
contenido que el de uno que ejecuta acciones directamente sobre una planta activa,
puesto que se limita a la lectura y análisis de datos históricos. Esto hace que el
despliegue del agente sea razonable con las tecnologías LLM actuales. Además, ese
carácter de solo lectura lo hace fácilmente escalable y transversal, pudiendo aplicarse
a otros ámbitos más allá del consumo eléctrico.

\section{Arquitectura de supervisión industrial}
\label{sec:arquitectura-supervision-industrial}

Antes de bajar al diseño concreto conviene tener clara la lógica que
organiza los sistemas de una planta: cómo se relacionan los sistemas de
negocio con los de control. La referencia habitual es la pirámide de
automatización, derivada del modelo de referencia de Purdue y formalizada
por el estándar ISA-95, que ordena la planta en cinco niveles, del proceso
físico en la base a la planificación empresarial en la cúspide
\cite{craci-arxiv}. La Figura~\ref{fig:piramide-isa95} recoge esta
jerarquía.

\begin{figure}[tphb]
  \centering
  \includegraphics[width=0.85\textwidth]{fig4_piramide_isa95}
  \caption{Jerarquía funcional de la automatización industrial según el modelo ISA-95.
    Elaboración propia a partir de \cite{craci-arxiv}.}
  \label{fig:piramide-isa95}
\end{figure}

En la base se sitúa el nivel 0, el proceso físico propiamente dicho:
sensores y actuadores, el hardware que mide y modifica las magnitudes reales
de la planta. Justo encima, el nivel 1 agrupa los elementos que controlan
esa operación, con \gls{PLC} o \gls{RTU}.

Los sistemas SCADA ocupan el nivel 2, y es la capa sobre la que pivota este
trabajo. Su función consiste en integrar la adquisición y transmisión de
datos con el software de interfaz de usuario, de forma que toda la
información de campo llegue centralizada al responsable de supervisión y
este pueda intervenir sobre el proceso en tiempo real \cite{nist-sp800-82}. Como la
plataforma empleada reproduce esta capa mediante módulos independientes,
conviene detenerse en sus componentes: el servidor de control, que almacena
y procesa lo que llega de los PLC o RTU; la interfaz \gls{HMI}, donde se
dispone la información al operador y desde la que este puede interactuar
sobre el proceso; el historiador de datos, que guarda la evolución temporal
de las variables de proceso y sobre el que se apoya el análisis posterior;
y las estaciones de ingeniería, desde las que se configura y mantiene el
resto de elementos. Estos cuatro componentes suelen comunicarse entre sí
dentro de una misma red local; hacia la capa de dispositivos de campo, en
cambio, se recurre a protocolos específicos según la dispersión geográfica
o el volumen de dispositivos \cite{nist-sp800-82}.

Por encima quedan dos niveles que este trabajo no aborda directamente. El
nivel 3 corresponde a la gestión de operaciones de fabricación, con
sistemas \gls{MES} que coordinan la producción y llevan cuenta de la
trazabilidad y el mantenimiento; el nivel 4, a la planificación
empresarial, con sistemas \gls{ERP} que gestionan compras, ventas o
finanzas.

El desarrollo descrito en esta memoria se sitúa en el nivel 2. El
historiador se corresponde con la base de datos histórica del proyecto,
alimentada a partir de las variables que el SCADA recoge del proceso; la
interfaz de supervisión que se desarrolla —descrita en el
apartado~\ref{sec:plataforma-n3uron}— cumple el papel de la HMI; y el
backend de configuración cumple el del servidor de control. La conexión con
capas superiores queda fuera del alcance del proyecto: no se integra con un
MES ni con un ERP, aunque el diseño no impide una futura extensión en ese
sentido.

\section{La plataforma N3uron}
\label{sec:plataforma-n3uron}

Como se mencionó en el punto anterior, el nivel 2 de la pirámide de automatización es en
torno al cual gira el presente TFG. Para ello, se optó por emplear la plataforma N3uron
debido a varias características que la hacen óptima para cumplir los objetivos del
trabajo. Esta plataforma, de tipo IIoT/DataOps, se organiza en torno a un modelo común
de datos —los \emph{tags}— sobre el que operan los distintos módulos que ofrece
\cite{n3uron-platform-intro}. Para la realización de este trabajo, se emplearán los
siguientes de su catálogo.

WebUI, el módulo principal, encargado de orquestar, controlar y supervisar el resto de
módulos, tags o protocolos. En términos del modelo ISA-95, se corresponde con el papel
del servidor de control, alojando la lógica de supervisión y centralizando la
configuración del sistema.

Historian, que tal y como indica su nombre es el encargado de la historización de
variables, con soporte para varios motores según el volumen y el tipo de despliegue; en
este caso, basado en MongoDB \cite{n3uron-historian-intro}. El resto de módulos de la
plataforma tiene acceso a los datos que almacena, lo que lo consolida como la pieza
central del trabajo, actuando como el historiador de datos descrito en el
punto~\ref{sec:arquitectura-supervision-industrial}.

WebVision, sobre el cual se despliega la capa de visualización, con elementos
predeterminados como gráficas, o elementos interactivos y de control, que permiten
desplegar el HMI descrito en el punto~\ref{sec:arquitectura-supervision-industrial}
\cite{n3uron-webvision}.

Scripting, a partir del cual se puede implementar lógica personalizada mediante código
\cite{n3uron-scripting}. Se empleará para diversas tareas como la importación de los
históricos al sistema o la creación de herramientas personalizadas para el MCP.

MCP Server, que actúa como puente de información entre N3uron y los modelos de
lenguaje. N3uron ha incorporado esta tecnología de forma nativa en su plataforma,
poniendo a disposición del usuario documentación, herramientas y \emph{prompts}
predeterminados \cite{n3uron-mcp-server}.

\section{Modelos de lenguaje de gran tamaño}
\label{sec:modelos-lenguaje}

Los modelos de lenguaje de gran tamaño se basan en redes neuronales
entrenadas sobre grandes volúmenes de texto, con la finalidad de, dada una
entrada, devolver la continuación más probable. Aplicado de forma iterativa,
este mecanismo permite a los modelos mantener conversaciones y responder
preguntas sin haber sido programados específicamente para esas tareas. La
arquitectura adoptada por la gran mayoría de modelos es la denominada
\emph{transformer}, cuya característica distintiva es el mecanismo de
atención: la capacidad de ponderar, para cada palabra que genera, la
relevancia del resto de palabras del contexto, sin importar la distancia
entre ellas \cite{vaswani2017}.

Sin embargo, un LLM necesita algo más para ser útil en el ámbito concreto de
este trabajo: dos capacidades resultan determinantes para que N3uron, el
MCP y AnythingLLM puedan funcionar como un conjunto. La primera es la
ventana de contexto, la cantidad de texto que el modelo puede procesar en
una única solicitud, medida en tokens —fragmentos de palabra que
constituyen la unidad mínima de procesamiento—: cuanto más amplia sea, más
histórico puede incluirse en una consulta sin necesidad de resumirlo antes
\cite{google-longcontext}. Y hace falta, además, la capacidad de llamada a
herramientas (\emph{function calling}): que el modelo pueda solicitar la
ejecución de una función externa e incorporar su resultado a la respuesta,
que es en la práctica lo que el MCP pone a su disposición
\cite{google-functioncalling}.

El acceso a estos modelos se realiza mediante API en los proveedores en la
nube —Anthropic, OpenAI o Google son los más relevantes para este trabajo, y
los tres ofrecen ventanas de contexto amplias y soporte nativo para
llamadas a herramientas—: la aplicación no aloja el modelo, sino que envía
cada consulta al servidor del proveedor y recibe la respuesta a través de
internet. Para la arquitectura de este TFG esto tiene una implicación
directa: aunque el modelo resida en la nube, el resto del sistema —N3uron,
la base de datos histórica, la interfaz de supervisión— puede mantenerse
íntegramente en la instalación local, y lo único que sale de la red de la
planta es el contenido de cada consulta.

Junto a estos proveedores en la nube existen también alternativas de
despliegue local, con herramientas como Ollama, que permite ejecutar
modelos de código abierto —Llama, Qwen, Mistral, entre otros— directamente
sobre el propio hardware, sin tráfico de datos con proveedores externos
\cite{ollama-official}. En un escenario donde la privacidad sea la prioridad, esta es la
alternativa más llamativa para este trabajo, porque resuelve directamente
una de las preocupaciones identificadas en el
apartado~\ref{sec:barreras-adopcion}. Tiene, sin embargo, un inconveniente frente a las
opciones anteriores: al ejecutarse en un equipo local, el rendimiento queda
limitado por la capacidad de ese equipo, y un hardware modesto puede
traducirse en un desempeño insuficiente.

En este TFG, la conexión entre el modelo y el servidor MCP se realiza
mediante AnythingLLM \cite{anythingllm-docs}, que actúa como cliente MCP
frente al MCP Server de N3uron: es esta plataforma, y no el modelo, quien
gestiona esa conexión, con independencia de si el modelo corre en la nube o
en local. Esto tiene una consecuencia práctica importante: la elección
entre un modelo en la nube y uno local no condiciona la viabilidad de la
integración con N3uron, y además otorga libertad para elegir el proveedor
LLM, ya que todos quedan recogidos en la infraestructura de AnythingLLM. La
decisión entre unos y otros pasa entonces a depender de otros criterios:
coste, rendimiento y el grado de exposición de los datos que la
organización esté dispuesta a asumir.

\section{Model Context Protocol}
\label{sec:model-context-protocol}

El Model Context Protocol resuelve precisamente el problema planteado en el
apartado~\ref{sec:estado-arte}: cómo dar acceso estructurado a datos y
herramientas a un LLM sin tener que construir, para cada proyecto, un
mecanismo de integración distinto. Se trata de un estándar abierto,
introducido por Anthropic en noviembre de 2024 \cite{mcp-docs}.

Define una arquitectura cliente-servidor con tres roles. El host es la
aplicación de IA que interactúa con el usuario —en este TFG, AnythingLLM—.
Dentro de él se aloja el cliente, encargado de mantener la conexión con un
servidor concreto; en el otro extremo está el servidor, el programa externo
que expone sus datos al cliente —aquí, el módulo MCP de N3uron—. Cliente y
servidor se comunican mediante JSON-RPC 2.0, sobre dos transportes
posibles: \emph{stdio} para servidores locales que corren como proceso del propio
sistema, o HTTP con eventos enviados por el servidor (\emph{streamable HTTP}) para
servidores remotos accesibles por red \cite{mcp-architecture}.

Sus capacidades se organizan en tres primitivas. Los \emph{tools} son
funciones externas ejecutables que el modelo puede invocar, cada una
definida por un nombre, una descripción y un esquema JSON de los parámetros
que acepta y devuelve. Los \emph{resources} son fuentes de datos de solo
lectura, identificadas mediante una URI, que el modelo puede consultar sin
que ello implique efectos secundarios sobre el sistema. Y los
\emph{prompts} son plantillas de instrucciones reutilizables, gestionadas
por el servidor, que guían al modelo hacia el uso más adecuado de
herramientas y recursos.

El flujo de una interacción sigue siempre el mismo patrón, recogido en la
Figura~\ref{fig:flujo-mcp}. El cliente descubre primero qué herramientas, recursos y
prompts ofrece el servidor. Cuando el modelo decide que necesita ejecutar
una acción, el host dirige al cliente para que envíe la solicitud
correspondiente; el servidor ejecuta entonces la lógica asociada —en este
trabajo, por ejemplo, una consulta al Historian de N3uron— y devuelve el
resultado, que el cliente incorpora al contexto del modelo. Con ese dato ya
disponible, el modelo genera su respuesta final \cite{mcp-architecture}.

\begin{figure}[tphb]
  \centering
  \includegraphics[width=0.85\textwidth]{fig5_flujo_mcp}
  \caption{Flujo de una interacción vía Model Context Protocol, desde la consulta del
    usuario hasta la respuesta informada por el dato histórico. Elaboración propia a
    partir de \cite{mcp-architecture}.}
  \label{fig:flujo-mcp}
\end{figure}

En este trabajo, el protocolo es lo que traduce el histórico de consumo
eléctrico almacenado en el Historian (apartado~\ref{sec:arquitectura-supervision-industrial}) en un
recurso consultable por el agente LLM, y las operaciones de análisis en
herramientas que ese agente puede invocar de forma autónoma a través de
AnythingLLM.

%%% Local Variables:
%%% TeX-master: "../book"
%%% End:
